% =====================================================================
%  CHAPTER 3 -- METHODOLOGY / REQUIREMENTS ANALYSIS
%
%  RUBRIC NOTE (R4): the top band requires the methods, approaches,
%  tools, techniques and algorithms to be "sufficiently discussed with
%  sufficient details and supporting figures". Two things follow:
%    (a) write it so another student could REPRODUCE the work;
%    (b) include at least one figure in this chapter (a process model
%        diagram is the obvious candidate -- a slot is reserved below).
% =====================================================================
\chapter{METHODOLOGY}
\label{ch:methodology}

This chapter presents the approach taken to build \fypProjectTitle{}:
the software process model followed, the requirements gathered and
analysed, and the tools, technologies and techniques used to implement
the system.

% ---------------------------------------------------------------------
\section{DEVELOPMENT APPROACH}
\label{sec:dev-approach}
% ---------------------------------------------------------------------
%  CARRIED OVER FROM PROPOSAL SECTIONS 4.1 and 4.3.
% ---------------------------------------------------------------------

This project follows the \textbf{Incremental Model}, in which the system
is developed and delivered in multiple functional increments. Each
increment adds features to, and improves upon, the previous version.
This approach permits early testing, easier debugging, and faster
delivery of working modules \cite{pan2021roles}.

The model suits this project specifically because the AI module carries
the highest technical risk and the least certain requirements. Building
a core functional system first -- registration, login and basic design
browsing -- allowed the platform to be tested end to end before the
speech recognition and semantic matching components were introduced, so
that failures in those components could be isolated from failures in the
surrounding application.

% EXPAND: add two or three sentences on how the increments actually ran
% in practice, including anything that had to be resequenced. Examiners
% value an honest account of a plan that changed over a tidy account of
% one that did not.

\subsection{Incremental development phases}
\label{subsec:increments}

\begin{enumerate}[label=\textbf{Increment \arabic*:}, itemsep=2pt,
                  labelwidth=2.8cm, labelsep=0.4em, align=left,
                  leftmargin=3.2cm]
  \item User registration, login and portfolio viewing.
  \item Design upload and management by architects, with semantic
        tagging.
  \item Construction company integration and collaboration features,
        real-time communication, and admin controls with verification.
  \item AI-based semantic and voice search functionality, and payment
        integration.
  \item Final testing, optimisation and deployment.
\end{enumerate}

\subsection{Major modules}
\label{subsec:major-modules}

% CARRIED OVER FROM PROPOSAL SECTION 4.2.
The system is decomposed into five modules, which are carried forward
into the subsystem architecture in Section~\ref{subsec:arch-subsystem}:

\begin{itemize}[leftmargin=*, itemsep=2pt]
  \item \textbf{User Module} -- user registration, authentication,
        browsing of architectural designs, and initiating contact with
        architects.
  \item \textbf{Architect Module} -- design uploads, portfolio
        management, and direct communication with clients using semantic
        tagging.
  \item \textbf{Construction Module} -- allows construction companies to
        list services and connect with clients through verified
        credentials.
  \item \textbf{Admin Module} -- administrative control for user
        verification, data moderation and analytics.
  \item \textbf{AI Module} -- a separate FastAPI service that embeds
        project briefs and architect portfolios using a Sentence-BERT
        model and ranks architects by cosine similarity
        \cite{reimers2019sbert}.
\end{itemize}

% =====================================================================
%  !!! DISCREPANCY BETWEEN PROPOSAL AND DELIVERED SYSTEM !!!
%  DECIDE WHAT TO DO ABOUT THIS BEFORE YOU SUBMIT.
%
%  Your proposal states that the AI Module integrates THREE components:
%      WhisperASR  -- speech-to-text for voice queries
%      SpaCy       -- keyword and entity extraction
%      S-BERT      -- semantic search
%
%  Only S-BERT is actually implemented. As of this review of the
%  repository:
%    - nlp-service/requirements.txt contains fastapi, uvicorn,
%      sentence-transformers, scikit-learn, httpx, python-dotenv and
%      pydantic. There is no whisper package and no spacy package.
%    - Searching the whole codebase (backend app/ and routes/, frontend
%      src/, nlp-service/) for "whisper" or "spacy" returns NO matches.
%    - nlp-service/matcher.py implements S-BERT embedding plus cosine
%      similarity only.
%
%  YOU HAVE THREE HONEST OPTIONS:
%
%  (a) IMPLEMENT THEM before submission, then this section is correct
%      as your proposal wrote it. Whisper is the larger job; SpaCy
%      keyword extraction is comparatively small.
%
%  (b) DESCRIBE THE SYSTEM AS BUILT -- which is what this report now
%      does -- and move voice input and keyword extraction into
%      Section~\ref{sec:future-work} as planned work. Then say so
%      plainly in the scope section and in the conclusion. This is
%      completely acceptable; FYP scope changes between proposal and
%      delivery all the time, and examiners respect it being stated.
%
%  (c) If a voice or SpaCy feature EXISTS but lives outside this
%      repository, add it to the repo or document where it lives.
%
%  WHAT YOU MUST NOT DO is leave the report claiming Whisper and SpaCy
%  are part of the delivered system while the code says otherwise. That
%  is the single most dangerous thing in this document: a demo or viva
%  question about voice input would expose it immediately.
%
%  Every place this affects is marked "DISCREPANCY" -- search for it.
% =====================================================================

% ---- Optional process-model figure (recommended for R4) --------------
\begin{figure}[H]
  \centering
  \fypdiagram[6cm]{[INSERT PROCESS MODEL DIAGRAM HERE]}
  \caption{Development process model followed in this project}
  \label{fig:process-model}
\end{figure}

% ---------------------------------------------------------------------
\section{REQUIREMENTS ANALYSIS}
\label{sec:requirements-analysis}
% ---------------------------------------------------------------------

\subsection{Requirements elicitation}
\label{subsec:elicitation}

% TODO: say HOW you gathered requirements -- interviews, questionnaires,
% observation, document analysis, competitor study. Give numbers where
% you have them (how many people, over what period).
[ELICITATION: the techniques used to gather requirements, who was
consulted, and over what period. Quantify wherever possible.]

\subsection{Functional requirements}
\label{subsec:functional-requirements}

% DERIVED FROM THE DELIVERED SYSTEM -- these are reverse-engineered
% from the 114 endpoints in archintent-backend/routes/api.php, grouped
% by module. Every row below corresponds to shipped functionality, so
% you can defend each one in the viva by pointing at a route.
%
% TODO: check the Priority column. It is set to a sensible default based
% on whether the feature is core to the main workflow, but priority is
% your judgement, not something the code can tell us.
{\small
\begin{longtable}{|l|p{9.2cm}|l|}
  \caption{Functional requirements of the system}
  \label{tab:functional-requirements} \\
  \hline
  \textbf{ID} & \textbf{Requirement} & \textbf{Priority} \\
  \hline
  \endfirsthead
  \hline
  \textbf{ID} & \textbf{Requirement} & \textbf{Priority} \\
  \hline
  \endhead
  \hline
  \endfoot

  \multicolumn{3}{|l|}{\textbf{Authentication and account management}} \\
  \hline
  FR-01 & The system shall allow a user to register as a client,
          architect or contractor. & High \\
  \hline
  FR-02 & The system shall verify a user's email address by one-time
          password before granting full access. & High \\
  \hline
  FR-03 & The system shall verify a user's phone number by one-time
          password. & Medium \\
  \hline
  FR-04 & The system shall authenticate users by email and password and
          issue an API token. & High \\
  \hline
  FR-05 & The system shall allow users to view and update their profile,
          including a profile image. & Medium \\
  \hline

  \multicolumn{3}{|l|}{\textbf{Client and project management}} \\
  \hline
  FR-06 & The system shall allow a client to create a project specifying
          project type, location and a written brief. & High \\
  \hline
  FR-07 & The system shall allow a client to view, update and delete
          their own projects. & High \\
  \hline
  FR-08 & The system shall generate a ranked list of matching architects
          for a project brief. & High \\
  \hline
  FR-09 & The system shall allow a client to select an architect for a
          project. & High \\
  \hline

  \multicolumn{3}{|l|}{\textbf{Architect module}} \\
  \hline
  FR-10 & The system shall allow an architect to create a portfolio and
          add portfolio projects with images and style tags. & High \\
  \hline
  FR-11 & The system shall allow an architect to set a cover image for
          and delete a portfolio project. & Medium \\
  \hline
  FR-12 & The system shall allow an architect to deliver a design and to
          receive revision requests against it. & High \\
  \hline
  FR-13 & The system shall allow an architect to onboard to Stripe
          Connect to receive payouts. & High \\
  \hline

  \multicolumn{3}{|l|}{\textbf{Contractor module}} \\
  \hline
  FR-14 & The system shall allow a contractor to create a portfolio with
          projects and images. & High \\
  \hline
  FR-15 & The system shall allow a client to post a project for
          construction bidding. & High \\
  \hline
  FR-16 & The system shall allow a contractor to submit a bid against a
          posted construction job. & High \\
  \hline
  FR-17 & The system shall allow a client to accept or reject a bid. & High \\
  \hline

  \multicolumn{3}{|l|}{\textbf{Agreements and payments}} \\
  \hline
  FR-18 & The system shall generate a digital agreement for a project. & High \\
  \hline
  FR-19 & The system shall allow parties to sign and finalise an
          agreement. & High \\
  \hline
  FR-20 & The system shall process project payments through Stripe and
          record them. & High \\
  \hline
  FR-21 & The system shall support refunding a recorded payment. & Medium \\
  \hline
  FR-22 & The system shall maintain a credit wallet, allow purchase of
          credit packages, and record credit transactions. & Medium \\
  \hline

  \multicolumn{3}{|l|}{\textbf{Communication}} \\
  \hline
  FR-23 & The system shall allow users to start a conversation and
          exchange messages in real time. & High \\
  \hline
  FR-24 & The system shall report unread message counts and allow a
          conversation to be marked as read. & Medium \\
  \hline

  \multicolumn{3}{|l|}{\textbf{Reviews}} \\
  \hline
  FR-25 & The system shall allow a client to review an architect or
          contractor after an eligible project. & Medium \\
  \hline
  FR-26 & The system shall display reviews on architect and contractor
          profiles. & Medium \\
  \hline

  \multicolumn{3}{|l|}{\textbf{Administration}} \\
  \hline
  FR-27 & The system shall allow an administrator to review the identity
          documents of pending architects and contractors and to verify
          or reject them. & High \\
  \hline
  FR-28 & The system shall allow an administrator to suspend, activate
          or delete a user account. & High \\
  \hline
  FR-29 & The system shall record administrative actions in an audit
          log. & Medium \\
  \hline
  FR-30 & The system shall present platform analytics on an
          administrative dashboard. & Low \\
  \hline
\end{longtable}
}

\subsection{Non-functional requirements}
\label{subsec:nonfunctional-requirements}

% Rows marked (implemented) were read out of the codebase and are
% defensible today. Rows marked [TARGET] are the ones only you can
% set -- they are the measurable goals Chapter 6 has to report against,
% so fill them in BEFORE you run your evaluation, not after.
\begin{table}[H]
  \centering
  \caption{Non-functional requirements of the system}
  \label{tab:nonfunctional-requirements}
  \small
  \begin{tabularx}{\textwidth}{|l|l|X|}
    \hline
    \textbf{ID} & \textbf{Attribute} & \textbf{Requirement} \\
    \hline
    NFR-01 & Security &
      The system shall authenticate API requests using Laravel Sanctum
      bearer tokens, and shall store passwords only as hashes
      (implemented). \\
    \hline
    NFR-02 & Security &
      The system shall rate-limit authentication endpoints to mitigate
      brute-force attacks: 5 login attempts, 6 registrations and 6
      one-time password requests per minute per client
      (implemented). \\
    \hline
    NFR-03 & Security &
      The system shall enforce role-based authorisation across the four
      roles client, architect, contractor and administrator
      (implemented). \\
    \hline
    NFR-04 & Security &
      The system shall not activate an architect or contractor account
      for public listing until an administrator has reviewed the
      submitted identity document (implemented). \\
    \hline
    NFR-05 & Auditability &
      The system shall record administrative actions in an audit log,
      and shall soft-delete user records rather than removing them
      (implemented). \\
    \hline
    NFR-06 & Performance &
      A semantic match request shall return a ranked architect list
      within [TARGET, e.g. 3] seconds for a catalogue of
      [TARGET, e.g. 500] portfolio projects. \\
    \hline
    NFR-07 & Performance &
      Chat messages shall be delivered to a connected recipient within
      [TARGET, e.g. 1] second over the WebSocket connection. \\
    \hline
    NFR-08 & Usability &
      The interface shall be responsive and usable on viewport widths
      from [TARGET, e.g. 360] px upward, without horizontal
      scrolling. \\
    \hline
    NFR-09 & Reliability &
      The platform shall remain available when the NLP service is
      unreachable, degrading search rather than failing the request.
      [VERIFY: confirm this is the actual behaviour, then keep or drop
      this row.] \\
    \hline
    NFR-10 & Maintainability &
      The NLP capability shall be deployable independently of the main
      application (implemented as a separate FastAPI service). \\
    \hline
  \end{tabularx}
\end{table}

\subsection{Feasibility study}
\label{subsec:feasibility}

% TODO: keep this short -- one paragraph per feasibility type.
[TECHNICAL FEASIBILITY: were the required technologies available and
within the team's competence?]

[OPERATIONAL FEASIBILITY: would the intended users actually adopt the
system in their existing workflow?]

[ECONOMIC FEASIBILITY: the cost of building and running the system
against the benefit it delivers.]

% ---------------------------------------------------------------------
\section{TOOLS AND TECHNOLOGIES}
\label{sec:tools}
% ---------------------------------------------------------------------

% VERSIONS READ FROM THE REPOSITORY, not from the proposal:
%   archintent-backend/composer.json + composer.lock  (locked versions)
%   archintent-frontend/package.json                  (semver ranges)
%   nlp-service/requirements.txt                      (pinned ==)
%
% Backend versions are the LOCKED ones from composer.lock, which is what
% actually ran. Frontend entries are the semver ranges from
% package.json -- if you want the exact resolved versions, read them out
% of package-lock.json and replace the carets.
\begin{table}[H]
  \centering
  \caption{Tools and technologies used in the project}
  \label{tab:tools}
  \small
  \begin{tabularx}{\textwidth}{|l|l|l|X|}
    \hline
    \textbf{Category} & \textbf{Technology} & \textbf{Version} &
    \textbf{Purpose} \\
    \hline
    Language & PHP & 8.2 & Backend runtime \\
    \hline
    Backend & Laravel & 12.57.0 & Application logic and REST API \\
    \hline
    Backend & Laravel Sanctum & 4.3.1 & Token-based API authentication \\
    \hline
    Backend & Laravel Reverb & 1.10.0 & WebSocket server for real-time chat \\
    \hline
    Database & MySQL & [VERSION] & Persistent relational storage \\
    \hline
    Frontend & React & 18.x & Component-based user interface \\
    \hline
    Frontend & TypeScript & 4.x & Static typing for the frontend \\
    \hline
    Frontend & Vite & 8.x & Build tool and development server \\
    \hline
    Frontend & Tailwind CSS & 3.x & Styling and responsive layout \\
    \hline
    Frontend & React Router & 6.x & Client-side routing \\
    \hline
    Frontend & Axios & 0.21.1 & HTTP client for API calls \\
    \hline
    Frontend & Laravel Echo / Pusher-js & 2.3.4 / 8.5.0 & Real-time event subscription \\
    \hline
    Frontend & Recharts & 3.8.1 & Dashboard data visualisation \\
    \hline
    AI service & Python + FastAPI & 0.115.0 & NLP microservice \\
    \hline
    AI service & sentence-transformers & 3.3.1 & S-BERT embeddings \\
    \hline
    AI service & scikit-learn & 1.5.2 & Cosine similarity scoring \\
    \hline
    AI service & Uvicorn & 0.30.6 & ASGI server for the NLP service \\
    \hline
    Payments & stripe-php & 20.1.0 & Server-side payment processing \\
    \hline
    Payments & @stripe/react-stripe-js & 6.1.0 & Client-side payment elements \\
    \hline
    Testing & PHPUnit & 11.5.3 & Backend unit and feature tests \\
    \hline
    Version control & Git and GitHub & --- & Source control and collaboration \\
    \hline
    Development & Visual Studio Code & --- & Primary IDE \\
    \hline
    Design & Figma & --- & UI/UX design and prototyping \\
    \hline
    Hosting & Vercel and Render & --- & Frontend and backend deployment \\
    \hline
  \end{tabularx}
\end{table}

% DISCREPANCY: WhisperASR and SpaCy appeared in the proposal's stack
% table but are NOT in any dependency file, so they are not listed here.
% See the discrepancy note in Section~\ref{subsec:major-modules}.
%
% TODO: MySQL version is the one value not recoverable from the repo --
% run  SELECT VERSION();  in your database and fill it in.

% TODO: justify the non-obvious choices. Your proposal lists the stack
% but does not defend it, and an unjustified stack reads as an accident
% rather than a decision. The three worth defending explicitly:
[JUSTIFICATION: explain (a) why Laravel was chosen for the backend over
the alternatives available to the team, (b) why S-BERT was chosen for
semantic matching rather than a keyword or TF-IDF approach -- Chapter 2
already supplies the argument and the citation, so refer back to
Section~\ref{subsec:lit-semantic} -- and (c) why the NLP components run
as a separate service rather than inside the Laravel application.]

% ---------------------------------------------------------------------
\section{TECHNIQUES AND ALGORITHMS}
\label{sec:techniques}
% ---------------------------------------------------------------------
%  DELETE THIS SECTION if your project does not involve a distinctive
%  algorithm or technique. Keep it if it does -- R4 names "algorithms"
%  explicitly, and a described algorithm is easy marks.
% ---------------------------------------------------------------------

% WRITTEN FROM THE ACTUAL IMPLEMENTATION in nlp-service/matcher.py.
% This section is where your project earns its "complex engineering
% problem" marks, so it is worth getting exactly right.

The matching engine is the central technique of this project. It is
implemented as a standalone FastAPI service that ranks architects
against a client's project brief by semantic similarity rather than by
keyword overlap \cite{reimers2019sbert}.

\subsection{Text representation}
\label{subsec:text-representation}

Two text representations are constructed. A \emph{project text} is
assembled from the project type, the location and the client's written
brief. An \emph{architect text} is assembled from the architect's
specialisation and biography, concatenated with the description and
style tags of one of their portfolio projects.

Crucially, one architect text is built for \emph{each} portfolio project
an architect owns, rather than one per architect. This lets the system
score an architect by their single best-matching work instead of by an
average over an uneven portfolio, or by whichever project happened to be
uploaded first. An architect with no portfolio projects is still
represented, by their base specialisation and biography text alone.

\subsection{Similarity scoring}
\label{subsec:similarity-scoring}

All texts -- the project text and every architect-project text -- are
encoded in a single batch using a Sentence-BERT model, by default
\texttt{all-MiniLM-L6-v2}, with the embeddings L2-normalised at encoding
time. Cosine similarity is then computed between the project vector and
every candidate vector.

% TODO: state the embedding dimensionality of the model you finally
% shipped with, and say why MiniLM was chosen over a larger model --
% the honest answer is usually the speed/accuracy trade-off, which is a
% perfectly good engineering justification.
[MODEL JUSTIFICATION: state the embedding dimensionality and why this
model was selected over larger alternatives.]

Each candidate score is collapsed to the best score per architect, then
converted to an integer in the range 0--100 by multiplying by one
hundred and rounding, clamped at both ends. Results below a configurable
minimum similarity, default 0.20, are discarded, and the top $k$
architects, default 10, are returned.

% TODO: both thresholds below are configuration you chose. Justify them,
% ideally with a sentence about what you observed when you tuned them.
[THRESHOLD JUSTIFICATION: explain how the minimum score of 0.20 and the
top-$k$ value of 10 were arrived at.]

\subsection{Integration with the main application}
\label{subsec:nlp-integration}

% TODO: describe the call path. The backend exposes an internal endpoint
% (POST /internal/project-matches) that receives the ranked results, so
% describe which side initiates, how the two services authenticate to
% each other, and whether the call is synchronous or deferred.
[INTEGRATION: describe how the Laravel application and the FastAPI
service exchange data, how they authenticate to one another, whether
matching happens synchronously on project creation or asynchronously,
and where the resulting matches are persisted -- the
\texttt{project\_matches} table.]

% DISCREPANCY -- SEE THE NOTE IN SECTION~\ref{subsec:major-modules}.
% The proposal describes this pipeline as three stages: Whisper
% speech-to-text, then SpaCy keyword extraction, then S-BERT matching.
% Only the third stage exists in the repository. This section documents
% what is implemented. If you build the other two before submission,
% add them back here as Sections 3.4.0 and 3.4.1 and restore the
% \cite{radford2022whisper} citation for the speech stage.

% EXPAND: a labelled diagram of this pipeline would strengthen R4
% considerably. A placeholder slot is available above at
% Figure~\ref{fig:process-model}.
